%% Document created 21 March 2021 automatically 
%% from /Users/massimosotgia/Desktop/uni_at_DIFI/Lab_C03/setup.py 

%% Copyright (C) Mattia Sotgia et al. 2021
%% Using class lab_unige.cls
%                                                            
%                                                            
%   **                 **             ******   ****   ****   
%  /**                /**            **////** *///** */// *  
%  /**        ******  /**           **    // /*  */*/    /*  
%  /**       ´´´´´´** /******      /**       /* * /*   ***   
%  /**        ******* /**///**     /**       /**  /*  /// *  
%  /**       **´´´´** /**  /** **  //**    **/*   /* *   /*  
%  /********//********/****** /**   //****** / **** / ****   
%  ////////  //////// /////   //     //////   ////   /´///   
%                                                            
%                                                            
\documentclass[italian, a4paper, 10pt, twocolumn]{../../style/lab_unige}
\usepackage[a4paper, margin=1.25cm, footskip=0.25in]{geometry}

\usepackage[utf8]{inputenc}
\usepackage[T1]{fontenc}

\usepackage[italian]{babel}

% \usepackage{biblatex}

\usepackage[bookmarksopen=true, 
citebordercolor={0 1 0}, 
linkbordercolor={1 0 0}, 
urlbordercolor={0 1 1}]{hyperref}
\usepackage[numbered]{bookmark}

\usepackage{graphicx}
\graphicspath{{../fig/}}
\usepackage{array}
\usepackage{tabulary}
\usepackage{booktabs}

% FOUNDAMENTAL
\usepackage{../../style/custom}

\usepackage{physics}

\usepackage{breqn}
\usepackage{cuted}
\usepackage{txfonts}

\usepackage{lipsum}

%% Define ref types
\newcommand{\reftab}[1]{Tab. {\ref{#1}}}%
\newcommand{\reffig}[1]{Fig. {\ref{#1}}}%
\newcommand{\refeqn}[1]{({\ref{#1}})}%
%% PAPER ONLY custom Macros
\newcommand{\Oi}{$O_1$}
\newcommand{\Oii}{$O_2$}
\newcommand{\treSigma}{$3\sigma$}
% \newcommand{\treSigma}{$3\sigma$\space}
\newcommand{\stdNG}[2]{$S_{#1}$($#2$)} %<- ???
\newcommand{\stdErr}[1]{$\varepsilon_{#1}$}
\newcommand{\mstdErr}[1]{\varepsilon_{#1}}
\newcommand{\gLab}{$g_t~=~(~9.8056~\pm~0.0001~\text{ stat}~)~\text{ m/s}^2$}
\newcommand{\ChiSqr}{$\chi^2$\space}
\newcommand{\ChiNdf}{$\chi^2/\text{ndf}$\space}
\newcommand{\cernroot}{\texttt{root} \space}
\newcommand{\lr}{$l_{\text{R}}$}
\newcommand{\Ti}[1]{$T_{#1}$}
\newcommand{\Ma}{$M_a$}
\newcommand{\Mb}{$M_b$}
\newcommand{\Tiso}{$T^*$}
\newcommand{\xiso}{$x^*$}
\newcommand{\LrVALUE}{$l_R=(0.800010\pm0.000010)$~mm}


%%
\setlength{\columnsep}{6mm}

\begin{document}
    \twocolumn[
    \begin{@twocolumnfalse}
        \title{
            {\raggedright \includegraphics[width=0.2\linewidth]{../../style/lab_mark.pdf}\\}
            Misura Sperimentale dell'Accelerazione di Gravità Sfruttando il Pendolo di Kater
        }
        \author{
        Eugenio Dormicchi\textsuperscript{1, 2},
        % Riccardo Pizzimbone\textsuperscript{1}, 
        Giovanni Oliveri\textsuperscript{1},
        Mattia Sotgia\textsuperscript{1}
        }

        \date{
        \textsuperscript{1}Gruppo C03, Esperienza di laboratorio n. 6 \\
        %\textsuperscript{2}In presenza in laboratorio per la presa dati\\
            % Università degli Studi di Genova, Dipartimento di Fisica.\\
            Presa dati-- 
            24 Marzo 2021, 15:00– 18:00; Analisi dati-- 
            30 Marzo 2021
        }
        \maketitle
        
        \begin{abstract}
            \textit{Obiettivo-- }
            Si vuole sfruttare uno strumento di precisione come il pendolo di Kater per ricavare dalla misura accuraata del periodo di oscillazione un valore atendibile della misura dell'accelerazione di gravità $g$.
            \textit{Metodi-- }
            Sfruttando il periodo di isocronia \Tiso~ del pendolo di Kater, ricavato dalla posizione $x_b$ della massa mobile \Mb , possiamo ricavare il valore di $g$.
            \textit{Risultati-- }
        
            \textit{Conclusione-- }
        
        
        \end{abstract}
        \vspace{2em}
    \end{@twocolumnfalse}
    ]

    %%%% CORPO DEL TESTO
    %%%% CORPO DEL TESTO

    \section{Obiettivo}
    \label{section:aim}
    Il valore della accelerazione di gravità $g$ può essere ottenuto in diversi metodi, con strumenti via via più precisi, e di conseguenza con risultati che si spera permettano una miglior approssimazione del valore vero di $g$. Il pendolo può essere sfruttato come metodo per misurare tale grandezza, e si è già eseguito l'esperimento sfruttando il pendolo semplice. Per questo motivo dal pendolo semplice si è passati a considerare il pendolo fisico, in particolar modo il pendolo di Kater, che permette di estendere la natura del pendolo 
    "semplice" e anche di eseguire misure con un livello maggiore di dettaglio.\\
    Sfruttando quindi una precisione maggiore e proprietà particolari del pendolo di Kater come pendolo fisico, vogliamo quindi individuare il valore del periodo di isocronia \Tiso~ caratteristico del pendolo di Kater preso in considerazione e da questo valore ricavare una buona stima di $g$, verificandone infine la compatibilità con il valore noto \gLab.

    \section{Strumentazione}
    \label{section:strument}
    Abbiamo a disposizione i seguenti strumenti:\\
    Calibro ventesimale, di portata 150 mm e sensibilità uguale all'accuratezza dello strumento di 50 $\mu$m;\\
    Pendolo di Kater (si veda in seguto descrizione dettagliata dell'apparato);\\
    Cronometro elettronico collegato a una fotocellula in grado di misurare una singola oscillazione del pendolo, di portata potenzialmente infinita (molto maggiore comunque dei periodi misurati), e sensibilità 10$^{-7}$ s, che assumiamo anche come accuratezza;\\

    Il pendolo di Kater è costituito da due masse \Ma~(massa rossa) e \Mb~(massa blu), poste su un'asta di acciaio inox., una fissa (\Ma) e una libera di muoversi, vincolata da una vite che la può bloccare all'asta in ben definite posizioni, che sono indicate da $N=25$ fori equidistanti 25~mm tra loro, eseguiti con strumentazione abbastanza precisa da poter trascurare l'incertezza associata a questa distanza. Il pendolo è poi poggiato su un sostegno tramite due perni che si incastrano perfettamente in una sede su una piastra posta orizzontalmente munita di livella a bolla, mantenuta in perfetto equilibrio.
    Il pendolo di Kater viene anche definito pendolo reversibile, in quanto possiamo agganciare il pendolo sia tramite i perni posti sull'astremo vicino alla massa \Ma, che su due perni identici posti sull'estremo della massa \Mb, per consentire quindi di misurare due periodi diversi \Ti{1} e \Ti{2}. 

    \section{Metodi}
    \label{section:methods}
    Tutte le misure sono riportate nelle unità del Sistema Internazionale (SI/MKS, ovvero metri, kilogrammi, secondi)
    
    %% MAGARI AGGIUNGERE COMMENTO SU QUANTE CIFRE SI CONSIDERANO E COME APPROSSIMIAMO I VALORI %%
    I valori riportati sono stati approssimati tenendo conto di alcune convenzioni prese. Si approssima l'errore a una cifra significativa se tale cifra è $\geqslant3$, altrimenti se tale cifra è 1 o 2 allora si considerano due cifre significative. Considerando quindi le posizioni decimali significative dell'errore si approssima per eccesso il valore numerico della grandezza.\\

    Abbiamo già accennato alla maggiore accuratezza nell'utilizzare il pendolo di Kater per misurare con maggiore precisione il valore di $g$. Infatti il pendolo di Kater può essere posto nella condizione di isocronia, dove il periodo misurato con il pendolo in un verso, viene uguale al periodo del caso in cui il pendolo sia fatto oscillare per l'altro estremo di oscillazione.


    \section{Risultati}
    \label{section:results}
    Si fa spesso riferimento anche alla regola del \treSigma, con la quale si vuole intendere la volontà di trasformare un errore di tipo massimo in errore statistico, ovvero calcolando $\mstdErr{x}~=~\Delta x/\sqrt{3}$.\\% e quindi considerando il valore vero con una probabilità statistica del \treSigma~$\approx99.73\$ di probabilità del dato vero.\\
    Nella trattazione che segue consideriamo come limite per definire un dato accettabile un valore che rientra in \treSigma~ (tre deviazioni standard) dal valore vero.

    
    \section{Conclusione}
    \label{section:conclusion}

    \subsection{Controlli}

    \subsection{Possibili errori sistematici}

    \appendix

    \setcounter{table}{0}
    \renewcommand{\thetable}{A\arabic{table}}

    \section{Dati completi}
    %% this file is for all the tables that go in the appendix as additional data
\begin{table*}[t!]
    \centering
    \caption{Dati grezzi dei periodi \Ti{1} e \Ti{2} misurati alle diverse lunghezze $x_b$.}
    \footnotesize
    \label{table:raw_data}
    \begin{tabular}{l*{6}{c}}
        \hline\hline\\[-1.5ex]
%                              & p1        & p2                 & p3                 & p4        & p5        & p6    
%       dati prima di approx = & 0.89865   & 0.8736499999999999 & 0.8486499999999999 & 0.82365   & 0.79865   & 0.77365
                               & \multicolumn{6}{c}{Misure di $x_b$ calcolate (m)}                                                        \\[+0.5ex] 
                               & \multicolumn{6}{c}{$x_b=\left[L_{0,~1} + L_{0,~2} + \frac{d}{2} + (N_{\text{fori}}\cdot0.025)\right]$ m (errore massimo $\pm1\times10^{-3}$ m)} 
                               \\[+0.5ex] \hline \\[-1.5ex]
        $x_b$                  & 0.899     & 0.874              & 0.849              & 0.824     & 0.799     & 0.774                      \\[+0.5ex] \hline \\[-1.5ex]
                               & \multicolumn{6}{c}{Periodi $T_1$ presi 10 volte per ogni $i$-esimo valore di $x_b$ (s)}                               \\[+0.5ex]
                               & \multicolumn{6}{c}{$T_{1_{(n \times 1)}}, \ldots, T_{1_{(n \times 6)}}$ (errore massimo $\pm10^{-7}$ s)} \\[+0.5ex] \hline \\[-1.5ex]
        $T_{1_{(1 \times i)}}$ & 1.8433475 & 1.8241344          & 1.8010792          & 1.7842967 & 1.7646188 & 1.7469742                  \\[+0.5ex]
        $T_{1_{(2 \times i)}}$ & 1.8438321 & 1.8250191          & 1.8014794          & 1.7831672 & 1.7654354 & 1.7470051                  \\[+0.5ex]
        $T_{1_{(3 \times i)}}$ & 1.8436974 & 1.8243224          & 1.8021671          & 1.7844837 & 1.7653484 & 1.7474456                  \\[+0.5ex]
        $T_{1_{(4 \times i)}}$ & 1.8441519 & 1.8246017          & 1.8026592          & 1.7838530 & 1.7646805 & 1.7473891                  \\[+0.5ex]
        $T_{1_{(5 \times i)}}$ & 1.8438134 & 1.8238088          & 1.8024890          & 1.7842300 & 1.7648372 & 1.7476352                  \\[+0.5ex]
        $T_{1_{(6 \times i)}}$ & 1.8440039 & 1.8240085          & 1.8030486          & 1.7839835 & 1.7650848 & 1.7468288                  \\[+0.5ex]
        $T_{1_{(7 \times i)}}$ & 1.8434289 & 1.8238919          & 1.8031204          & 1.7841418 & 1.7649879 & 1.7475854                  \\[+0.5ex]
        $T_{1_{(8 \times i)}}$ & 1.8437870 & 1.8239454          & 1.8048308          & 1.7838625 & 1.7647196 & 1.7475552                  \\[+0.5ex]
        $T_{1_{(9 \times i)}}$ & 1.8432011 & 1.8243763          & 1.8040768          & 1.7843554 & 1.7644894 & 1.7489131                  \\[+0.5ex]
        $T_{1_{(10\times i)}}$ & 1.8436028 & 1.8242312          & 1.8036150          & 1.7840838 & 1.7641298 & 1.7474017                  \\[+0.5ex] \hline \\[-1.5ex]
                               & \multicolumn{6}{c}{Periodi $T_2$ presi 10 volte per ogni $i$-esimo valore di $x_b$ (s)}                               \\[+0.5ex]
                               & \multicolumn{6}{c}{$T_{2_{(n \times 1)}}, \ldots, T_{2_{(n \times 6)}}$ (errore massimo $\pm10^{-7}$ s)} \\[+0.5ex] \hline \\[-1.5ex]
        $T_{2_{(1 \times i)}}$ & 1.8162160 & 1.8069800          & 1.7978242          & 1.7904099 & 1.7837065 & 1.7793795                  \\[+0.5ex]
        $T_{2_{(2 \times i)}}$ & 1.8162461 & 1.8074008          & 1.7982385          & 1.7919054 & 1.7839124 & 1.7790879                  \\[+0.5ex]
        $T_{2_{(3 \times i)}}$ & 1.8164208 & 1.8066094          & 1.7983828          & 1.7912305 & 1.7838475 & 1.7794034                  \\[+0.5ex]
        $T_{2_{(4 \times i)}}$ & 1.8162922 & 1.8068046          & 1.7984080          & 1.7911800 & 1.7840367 & 1.7794454                  \\[+0.5ex]
        $T_{2_{(5 \times i)}}$ & 1.8165907 & 1.8073671          & 1.7984383          & 1.7912973 & 1.7837570 & 1.7793693                  \\[+0.5ex]
        $T_{2_{(6 \times i)}}$ & 1.8161985 & 1.8071854          & 1.7986348          & 1.7911650 & 1.7840481 & 1.7794584                  \\[+0.5ex]
        $T_{2_{(7 \times i)}}$ & 1.8163087 & 1.8080697          & 1.7989685          & 1.7908471 & 1.7845240 & 1.7793767                  \\[+0.5ex]
        $T_{2_{(8 \times i)}}$ & 1.8162481 & 1.8068933          & 1.7985901          & 1.7912831 & 1.7840374 & 1.7792549                  \\[+0.5ex]
        $T_{2_{(9 \times i)}}$ & 1.8167033 & 1.8072342          & 1.7986484          & 1.7910748 & 1.7848659 & 1.7794661                  \\[+0.5ex]
        $T_{2_{(10\times i)}}$ & 1.8170554 & 1.8069937          & 1.7986683          & 1.7911061 & 1.7842935 & 1.7793790                  \\[+0.5ex]
        \hline
    \end{tabular}
\end{table*}

    
\end{document}
